% Nouveau EGA V (cote 25): archive cote 25, site page 169 (PDF page 170).
% typescript. Draft compared: none.
% AI-transcribed by Claude Opus 5.5 (claude-opus-5-5, adaptive thinking, effort high), Message Batches API, batch ega25b (msgbatch_01A7n5YuXubUPNK1zsepi8RG), 2026-11-13; not yet checked by eye.
\scanpage{25/169}
\mspage{\hand{\unsure{4}}}

image inverse de la section nulle du fibré vectoriel $V((\underline{\Omega}^m_{X/Y}\otimes\underline{O}_{V^o})^{\hand{\vee}}\otimes\bigwedge^m\underline{N}\otimes\underline{O}_{V^o})$
par la section correspondant à l'homomorphisme $(*)_{m}$ ). Comme le premier memb-%
re de $(*_m)$ est de type fini, il s'ensuit d'ailleurs que $V^1$ est de présantati%
on finie sur $V^o$ . De même, $V''$ se définit comme le préschéma des zéros de
\[ (*_{m-1}) \qquad \bigwedge^{m-1}\underline{N}\otimes\underline{O}_{V^1} \longrightarrow \underline{\Omega}^{m-1}_{X/Y}\otimes\underline{O}_{V^1} \]
déduit de $(*)$ en passant à la puissance extérieure m-1 ème et induisant sur
$V^1$ . Pour l'invariance da sous-préschéma $V^2\subset V^1-V''$ , que je n'ai pas vrai-%
ment vérifiée et que je ne peux donc que conjecturer, je pense qu'on doit
pou\xout{x}voir préciser ce point de la façon suivante. L'homomorphisme
\[ (*\,*^1) \qquad \underline{E}\otimes\underline{O}_{V^1}\longrightarrow\underline{\Omega}^1_{X/Y}\otimes\underline{O}_{V^1} \]
déduit de $(*^o)$ en \xout{\unsure{induisant sur V,}} en composant avec l'homomorphisme
\uline{surjectif} $\underline{E}\otimes\underline{O}_{V^o}\to\underline{N}$ et en \xout{\unsure{composant}} induisant sur $V^1$ , a même conoyau
$\underline{W}$ que
\[ (*^1) \qquad \underline{N}\otimes\underline{O}_{V^1}\to\underline{\Omega}^1_{X/Y}\otimes\underline{O}_{V^1} \quad , \]
lequel \inserted{$\underline{W}$} est donc localement libre sur $V^1-V''$ ,\hand{(}de rang d\deleted{im}+1 en un point
x de $V^1-V''$ en lequel X est de dim rel \hand{d} sur Y\hand{)}. Il s'ensuit que\hand{``}la formation
du noyau de $(*^1)$ \over{c}{C}ommute à tout changement de base $T\to V^1$\inserted{$-V''$}\hand{''}\over{.}{,} \over{s}{S}oit $\underline{k}$ le
noyau de $(*^1)$. Si $\underline{K}$ est le noyau de $(*\,*^1)$ , qui est donc inversible sur
$V^1-V''$ , on voit \xout{donc} que $\underline{k}$ est un quotient de $\underline{K}$ . D'autre part, on a défini
page 15 un homomorphisme
\[ \hand{\underline{M}} : \ \underline{K}\longrightarrow\underline{\mathrm{Sym}}^2(\underline{W}) \]
au dessus de $V^1-V''$ (il peut, être commode de désigner par $\underline{K}$,$\underline{W}$ les Modules
sur $V^1-V''$ exclusivement ...), \deleted{et} \over{l}{L}e point qu'on aurait envie de vérifier
alors est qu'en fait cet homomorphisme se factorise par le quotient $\underline{k}$ de $\underline{K}$:
\[ \hand{m} : \ \underline{k}\longrightarrow\underline{\mathrm{Sym}}^2(\underline{W}) \ \deleted{$*$} ; \]
et que ce dernier homomorphisme (entre deux Modules dont la définition ne
\xout{si ce point est vérifié, alors} dépend plus de u) est en fait indépendant de
u. Si ce point est vérifié (et \uline{s'il est vrai}, un calcul élémentaire doit le

\mspage{\hand{169}}
